\section*{Filtrado mediante $C_1$ y $C_2$}

En función a la figura \ref{fig:amplificador-voltaje} se determinaron la relación matemática para establecer las frecuencias de corte a 6 [Hz] en función a las funciones de transferencia de cada etapa considerando para ello un polo dependiente del capacitor y la resistencia de entrada, logrando así obtener la frecuencia de corte deseada.

La relación matemática descrita por el polo se conceptualiza en \ref{eq:get-capacitors}

\begin{equation}
	C_n = \frac{1}{2\pi fR_n}
	\label{eq:get-capacitors}
\end{equation}

Por tanto luego de aplicar este proceso, se obtuvieron una capacitancia de $C_1 = 5.78\mu F$ y $C_2 = 10\mu F$ respectivamente para cada etapa del amplificador, así mismo la etapa de ganancia queda determinada en función a la resistencias desde $R_1$ a $R_4$.

